
\begin{proposition}\label{cdp:st:C.1.3}
For every $Y$ and $D$ in Statement \ref{cdp:st:C.1.2} and every $L\in\Pic(Y)$,
\[
 h^2(Y,T_Y\otimes L)\geq
 \max\left\{0,1-\frac{b_3(Y)}2\right\}.
\]
In particular, $b_3(Y)=0$ gives $H^2(Y,T_Y\otimes L)\ne0$ for
every $L$. A threefold in this class possessing a twist with vanishing
second tangent cohomology has $b_3(Y)\geq2$.
\end{proposition}
\begin{proof}
On the dense open subset of each Picard component supplied by
Statement \ref{cdp:st:C.1.2}, Statement \ref{cdp:st:C.1.1} becomes
\[
 h^2(Y,T_Y\otimes L)-h^1(Y,T_Y\otimes L)
 =1-b_3(Y)/2.
\]
The asserted lower bound holds there. The jump locus for the lower
bound on $h^2$ is closed analytic in the explicit family above and
contains a dense subset of the component. It therefore contains the
entire component. This proves the bound for all line bundles,
including torsion bundles and the trivial bundle. The final assertion
is its direct numerical consequence.
\end{proof}

Statements \ref{cdp:st:C.1.1}--\ref{cdp:st:C.1.3} prove the stated index obstruction with rational Betti-number
vanishing. No integral trivialization of either Chern class is needed
for the index. The existence of the divisor enters in the endpoint
vanishing, and a positive multiple of its actual Cartier divisor
enters the Picard family. These statements supply an index
obstruction; they make no assertion that the failed degree-two
vanishing holds.


